% Part: lambda-calculus
% Chapter: church-rosser
% Section: definitions-and-properties

\documentclass[../../../include/open-logic-section]{subfiles}

\begin{document}

\olfileid[hi]{lam}{cr}{dap}

\olsection{परिभाषा और गुण}

इस अध्याय में हम चर्च--रॉसर गुण की संकल्पना और उसके कुछ सामान्य गुणों का
परिचय देते हैं।

\begin{defn}[चर्च--रॉसर गुण, CR] पदों पर कोई संबंध $\xredone$
  \emph{चर्च--रॉसर गुण} को तभी और केवल तभी संतुष्ट करता है जब $M \xredone
  P$ और $M \xredone Q$ होने पर कोई ऐसा~$N$ अस्तित्व में हो कि $P
  \xredone N$ और $Q \xredone N$।
\end{defn}

हम लैम्ब्डा कलन को संगणन के ऐसे मॉडल के रूप में देख सकते हैं जिसमें
प्रसामान्य रूप के पद ``मान'' और पदों पर अपचयनीयता-संबंध ``परिकलन नियम''
हैं। चर्च--रॉसर गुण कहता है कि जब किसी परिकलन को आगे बढ़ाने की एक से अधिक
विधियाँ हों, तब भी व्यंजक का केवल एक ही मान होता है।

प्रारंभिक बीजगणित से उदाहरण लें: $4 \times (1+2) + 3$ की गणना एक से अधिक
विधियों से हो सकती है। इसे या तो $4 \times 3+3$ में अपचयित किया जा सकता
है (यदि पहले $1+2$ को~$3$ में अपचयित करें), अथवा $4 \times 1+4 \times
2+3$ में (यदि पहले वितरणशीलता से $4 \times (1+2)$ का अपचयन करें)। किंतु
दोनों को आगे $12+3$ में अपचयित किया जा सकता है।

यदि हम $\xredone$ को $\beta$-अपचयन मानें, तो सरलता से देखते हैं कि
चर्च--रॉसर गुण का एक परिणाम यह है: यदि किसी पद का प्रसामान्य रूप हो, तो वह
अद्वितीय है। मान लें $M$ अपचयित होकर $P$ और $Q$ बन सकता है और दोनों
प्रसामान्य रूप हैं। चर्च--रॉसर गुण से कोई ऐसा $N$ है जिसमें $P$ और $Q$
दोनों अपचयित होते हैं। चूँकि परिकल्पना से $P$ और $Q$ प्रसामान्य रूप हैं,
इसलिए उनका $N$ में अपचयन केवल तुच्छ अपचयन हो सकता है; अर्थात् $P$, $Q$ और
$N$ एक ही हैं। इससे किसी पद के \emph{उस} प्रसामान्य रूप की बात करना उचित
है।

अतः लैम्ब्डा कलन को संगणन के मॉडल के रूप में देखने पर किसी पद का प्रसामान्य
रूप उस पद से आरंभ होने वाली संगणना का ``अंतिम परिणाम'' माना जा सकता है।
ऊपर के उपप्रमेय का अर्थ है कि किसी संगणना का अंतिम परिणाम, यदि हो, केवल एक
होता है; जैसे $4 \times (1+2)+3$ की गणना का केवल एक परिणाम~$15$ है।

\begin{thm} \ollabel{thm:str}
  यदि कोई संबंध $\xredone$ चर्च--रॉसर गुण संतुष्ट करता है और $\xred$ वह
  सबसे छोटा संक्रामक संबंध है जिसमें $\xredone$ समाहित है, तो $\xred$ भी
  चर्च--रॉसर गुण संतुष्ट करता है।
\end{thm}

\begin{proof}
  मान लें
  \begin{align*}
    M & \xredone P_1 \xredone \dots \xredone P_m \text{ और}\\
    M & \xredone Q_1 \xredone \dots \xredone Q_n.
  \end{align*}
  हम ऊँचाई $m + 1$ और चौड़ाई $n + 1$ वाले पदों का जालक $N$ बनाकर प्रमेय
  सिद्ध करेंगे। $i$-वीं पंक्ति और $j$-वें स्तंभ के पद को $N_{i,j}$ से
  दर्शाते हैं।
  
  हम $N$ को इस प्रकार बनाते हैं कि $N_{i,j} \xredone N_{i+1,j}$ और
  $N_{i,j} \xredone N_{i,j+1}$। इसकी परिभाषा है:
  \begin{align*}
    N_{0,0} &= M \\
    N_{i,0} &= P_i && \text{यदि } 1 \le i \le m \\
    N_{0,j} &= Q_j && \text{यदि } 1 \le j \le n \\
  \intertext{और अन्यथा:}
    N_{i,j} &= R
  \end{align*}
  जहाँ $R$ ऐसा पद है कि $N_{i-1,j} \xredone R$ और $N_{i,j-1} \xredone
  R$। $\xredone$ के चर्च--रॉसर गुण से ऐसा पद सदा अस्तित्व में है।

  अब $N_{m,0} \xredone \dots \xredone N_{m,n}$ और $N_{0,n} \xredone
  \dots \xredone N_{m,n}$। ध्यान दें कि $N_{m,0}$, ~$P$ और $N_{0,n}$,
  ~$Q$ है। $\xred$ की परिभाषा से प्रमेय सिद्ध होता है।
\end{proof}

\end{document}
